\documentclass{uofa-eng-assignment}
\usepackage{caption}
\usepackage{hyperref}
\usepackage{tcolorbox}
\usepackage{boxedminipage}
\hypersetup{
    colorlinks = false,
    linkbordercolor = {white},
}
\captionsetup[table]{name=Tabla, labelsep=period}
\captionsetup[figure]{name=Figura, labelsep=period}

\newcommand{\hpl}{\hyperlink}
\newcommand{\hpt}{\hypertarget}
\newcommand*{\name}{Joan Sebastian Ramirez Sanchez}
\newcommand*{\id}{2223948}
\newcommand*{\course}{Introducción a la teoría de control}
\newcommand*{\assignment}[1]{Tarea #1}
\newcommand*{\dated}{\today}
\newcommand{\tf}{\textbf}

\begin{document}

\maketitle
\section*{Vehículo Aéreo No Tripulado Cuadricóptero}

\textbf{Contexto y aplicación.} 
Los cuadricópteros son plataformas de 6 grados de libertad (6-DOF) con fuerte acoplamiento no lineal. Su control preciso es esencial en logística autónoma, inspección de infraestructura y movilidad aérea urbana. Las no linealidades surgen de las transformaciones de orientación mediante ángulos de Euler

\subsection*{Variables del sistema}

\begin{center}
\begin{tabular}{lll}
\hline
\textbf{Símbolo} & \textbf{Unidades} & \textbf{Descripción} \\
\hline
$\phi, \theta, \psi$ & rad & Ángulos de Euler (roll, pitch, yaw) \\
$\Omega_i$ & rad/s & Velocidad angular del rotor $i$ \\
$k_T$ & N$\cdot$s$^2$ & Constante de empuje \\
$k_D$ & N$\cdot$m$\cdot$s$^2$ & Constante de par aerodinámico \\
$I_x, I_y, I_z$ & kg$\cdot$m$^2$ & Momentos de inercia principales \\
$L$ & m & Distancia del centro al rotor \\
\hline
\end{tabular}
\end{center}

\subsection*{Ecuaciones no lineales del sistema}

\begin{tcolorbox}[colback=brown!10!white, colframe=brown!60!black, title=No linealidades principales]
\begin{align*}
    \ddot{z} &= \frac{T}{m} \cos\phi \cos\theta - g 
\quad \text{(producto de trigonom.)} \\
I_x \ddot{\phi} 
&= \dot{\theta}\dot{\psi}(I_y - I_z) + \tau_\phi
\end{align*}
\end{tcolorbox}



%%%%%%%%%%%%%%%%%%%%%%%%%%%%%%%%%%%%%%%%%%%%%%%%%%%%%%%%%%%%%%%%%%
%%%%%%%%%%%%%%%%%%%%%%%%%%%%%%%%%%%%%%%%%%%%%%%%%%%%%%%%%%%%%%%%%%
%%%%%%%%%%%%%%%%%%%%%%%%%%%%%%%%%%%%%%%%%%%%%%%%%%%%%%%%%%%%%%%%%%
%%%%%%%%%%%%%%%%%%%%%%%%%%%%%%%%%%%%%%%%%%%%%%%%%%%%%%%%%%%%%%%%%%
%%%%%%%%%%%%%%%%%%%%%%%%%%%%%%%%%%%%%%%%%%%%%%%%%%%%%%%%%%%%%%%%%%

\section{Deducción matemática del sistema}
Pendiente 


%%%%%%%%%%%%%%%%%%%%%%%%%%%%%%%%%%%%%%%%%%%%%%%%%%%%%%%%%%%%%%%%%%
%%%%%%%%%%%%%%%%%%%%%%%%%%%%%%%%%%%%%%%%%%%%%%%%%%%%%%%%%%%%%%%%%%
%%%%%%%%%%%%%%%%%%%%%%%%%%%%%%%%%%%%%%%%%%%%%%%%%%%%%%%%%%%%%%%%%%
%%%%%%%%%%%%%%%%%%%%%%%%%%%%%%%%%%%%%%%%%%%%%%%%%%%%%%%%%%%%%%%%%%


\section{Punto de equilibrio}
Para hallar el punto de equilibrio del sistema, buscamos las condiciones de vuelo estacionario (hovering). 
Esto implica que el vehículo se encuentra en una posición fija, sin velocidades ni aceleraciones lineales o angulares
por lo tanto debemos tener
\begin{itemize}
    \item Aceleraciones lineales nulas: $\ddot{z}=\ddot{\phi}=\ddot{\theta}=\ddot{\psi}=0$
    \item velocidades lineales nulas: $\dot{z}=\dot{\phi}=\dot{\theta}=\dot{\psi}=0$
\end{itemize}
Además que necesitamos que la orientación del "dron" sea tal que no este desbalanceado entonces 
necesitamos que $\phi=\theta=0$ dado que son los ángulos de roll y pitch respectivamente, 
entonces el único ángulo que puede ser diferente de cero es el yaw $\psi$ 
pero dado que no afecta a la dinámica del sistema entonces podemos asumir que $\psi=0$ 
sin perdida de generalidad, por lo tanto el punto de equilibrio del sistema 
es (\textbf{explicar que es roll y pitch en la sección 1})
por lo tanto haciendo el analizis de las ecuaciones diferenciales del sistema entonces tenemos que el punto de equilibrio es
\begin{align*}
    0= \frac{T}{m} \cos 0 \cos 0 - g \implies T=mg
\end{align*}
esto para la ecuación diferencial de la aceleración lineal en el eje z, 
y para las ecuaciones diferenciales de las aceleraciones angulares entonces tenemos que
\begin{align*}
    0 = \dot{\theta}\dot{\psi}(I_y - I_z) + \tau_\phi \implies \tau_\phi=0 \\
\end{align*}
%%%%%%%%%%%%%%%%%%%%%%%%%%%%%%%%%%%%%%%%%%%%%%%%%%%%%%%%%%%%%%%%%%
%%%%%%%%%%%%%%%%%%%%%%%%%%%%%%%%%%%%%%%%%%%%%%%%%%%%%%%%%%%%%%%%%%
%%%%%%%%%%%%%%%%%%%%%%%%%%%%%%%%%%%%%%%%%%%%%%%%%%%%%%%%%%%%%%%%%%
%%%%%%%%%%%%%%%%%%%%%%%%%%%%%%%%%%%%%%%%%%%%%%%%%%%%%%%%%%%%%%%%##





\section{Linealización del sistema}
Dado que hasta el momento hemos trabajado con sistema de ecuaciones diferenciales de primer orden entonces vamos
el grado de la siguiente ecuación diferencial de segundo orden a primer orden para eso vamos a introducir nuevas variables 
con una condición dado que $T$ es fijo vamos a asumir que $T=1$ sin perdida de generalidad, entonces tenemos que

\begin{align*}
    \begin{cases}
        x_1=z, \quad x_2=\dot{z} \\
        x_3=\phi, \quad x_4=\dot{\phi} \\
        x_5=\theta, \quad x_6=\dot{\theta} \\
        x_7=\psi, \quad x_8=\dot{\psi}
    \end{cases}    
\end{align*}
Por lo tanto nos queda el siguiente sistema de ecuaciones diferenciales de primer orden
\begin{align*}
    \begin{cases}
        \dot{x}_1 = x_2 \\
        \dot{x}_2 = \frac{\tau_z}{m} \cos x_3 \cos x_5 - g \\
        \dot{x}_3 = x_4 \\
        \dot{x}_4 = \frac{1}{I_x} \left( x_6 x_8 (I_y - I_z) + \tau_\phi \right) \\
        \dot{x}_5 = x_6 \\
        \dot{x}_6 = \frac{1}{I_y} \left( x_4 x_8 (I_z - I_x) + \tau_\theta \right) \\
        \dot{x}_7 = x_8 \\
        \dot{x}_8 = \frac{1}{I_z} \left( x_4 x_6 (I_x - I_y) + \tau_\psi \right)
    \end{cases}
\end{align*}
donde el control de entrada del sistema es el vector $\tau=[\tau_z, \tau_\phi, \tau_\theta, \tau_\psi]$ que es el empuje total generado por los cuatro rotores, dado que asumimos que $T=1$ entonces el sistema queda completamente definido por el vector de torques $\tau$ y las variables de estado $x$.
Ahora lo que buscaremos es linealizar el sistema de ecuaciones diferenciales haciendo una expansión de Taylor 
alrededor del punto de equilibrio que encontramos en la sección anterior, entonces tenemos que 
el punto de equilibrio respecto a las variables de estado es $x_e=(z_0,0,0,0,0,0,\psi_0,0)$, empecemos con $\dot{x}_2$
entonces tenemos que
\begin{align*}
    x_2(x) = \frac{T}{m} \cos x_3 \cos x_5 - g &\approx f_2(x_e) + \nabla f_2(x_e) \cdot (x - x_e) \\
    &\approx f_2(x_e) + \frac{\partial f_2}{\partial x_3}\bigg|_{x_e} (x_3 -0) + \frac{\partial f_2}{\partial x_5}\bigg|_{x_e} (x_5 - 0) \\
    &\approx 0 - \frac{T}{m} \sin 0 \cos 0 (0) - \frac{T}{m} \cos 0 \sin 0 (0)  =0
\end{align*}
para $\dot{x}_4$ entonces tenemos que
\begin{align*}
    x_4(x) = \frac{1}{I_x} \left( x_6 x_8 (I_y - I_z) + \tau_\phi \right) &\approx f_4(x_e) + \nabla f_4(x_e) \cdot (x - x_e) \\
    &\approx f_4(x_e) + \frac{\partial f_4}{\partial x_6}\bigg|_{x_e} (x_6 - 0) + \frac{\partial f_4}{\partial x_8}\bigg|_{x_e} (x_8 - 0) + \frac{\partial f_4}{\partial \tau_\phi}\bigg|_{x_e} (\tau_\phi - 0) \\
    &\approx 0 + \frac{1}{I_x} (0) (I_y - I_z) (0 - 0) + \frac{1}{I_x} 0 (I_y - I_z) (0 - 0) + \frac{1}{I_x} (\tau_\phi - 0) = \frac{1}{I_x} \tau_\phi
\end{align*}
\end{document}
